\section{Related work}

The budget-indexed oracle value connects treatment allocation under a capacity constraint with inference for optimized value functions. \citet[Assumptions 4.1.1--4.1.3 and Theorem 4.1.1]{FengHongNekipelov2026} derive an asymptotic value-process law for binary constrained allocation using differentiability of the sorting operator, a regular decision boundary, and joint process convergence. Their conclusion concerns the limiting distribution of a plug-in process in that regular regime. Here, for a known \(d\)-cell alphabet and fixed overlap, the estimand is the same type of oracle capacity curve, but the conclusion is a finite-sample minimax bound for expected squared supremum loss over every symmetric budget interval. When \(d\) is fixed the risk is of order \(1/n\); when \(d\) grows it is of order \(\min\{1,d/[n\log(ed)]\}\). The causal law class permits positive-mass treatment-effect ties at decision thresholds. \citet[Section~4]{FirpoGalvaoParker2023} develop uniform tests for marginally optimized value functions using directional derivatives and resampling. Their optimization framework is more general, while the present result gives a rate for the particular causal capacity curve under fixed overlap. \citet[Sections~3.1 and~5]{SverdrupWuAtheyWager2025} evaluate targeting across budgets with multiple treatment arms using Qini curves and asymptotic pointwise confidence intervals at specified budgets. Our target optimizes the population policy separately at each budget and our loss compares an estimated curve with that oracle curve.

\citet[Section~5]{BhattacharyaDupas2012} study welfare from a data-based binary allocation under a budget constraint and give asymptotic confidence intervals for that policy's welfare. \citet[Sections~3 and~5]{LuedtkeVanDerLaan2016} analyze differentiability and inference for the mean outcome under an optimal individualized rule, including nonunique optima. \citet[Introduction]{Hirano2009} develop asymptotic comparisons of statistical treatment rules. These results frame allocation and optimal-value inference; here the estimand is the population oracle optimum at every budget in \(I=[b_0,1-b_0]\), and the guarantee is an all-procedure finite-sample minimax bound for expected squared supremum loss under fixed overlap and a known discrete alphabet. Policy-learning work instead assesses the welfare of a selected rule \citep{KitagawaTetenov2018,AtheyWager2021,Pellatt2023}. The unrestricted discrete scalar result \citep{CausalSmithUnrestricted2026} studies the \(b=1\) coordinate under the same causal class and has the same minimax order \(\min\{1,d/[n\log(ed)]\}\). The factor-one dual contraction in \cref{obj:prop:dual-representation} shows that estimating every coordinate in \(I\) attains that order simultaneously, including when \(I=[0,1]\).

The statistical tools also connect to estimation of functionals on large discrete alphabets. Polynomial approximation and factorial-moment statistics support estimation of nonsmooth functionals \citep{JiaoVenkatHanWeissman2015}, while moment matching supplies lower bounds for such problems \citep{CaiLow2011,JiaoHanWeissman2018}. Here, polynomial approximation is applied to cell contributions across shadow prices, and paired-cell moment matching captures the difficulty of estimating the capacity-indexed value.
